Attempt.  I want \(m^{2028}\equiv 1\pmod {2197}\) and I interpret the first occurrence as forcing the smallest such residue.

I factor the modulus as
\[
2197=11\cdot 199.
\]
Then the congruence should be checked modulo \(11\) and modulo \(199\).

Modulo \(11\), the possible orders divide \(10\).  In particular \(m\equiv 1\) and \(m\equiv 10\) both give \(m^{2028}\equiv 1\pmod {11}\).  Modulo \(199\), the orders divide
\[
\phi(199)=198.
\]
The simplest way to satisfy both congruences is to take
\[
m\equiv 1\pmod {11},\qquad m\equiv 1\pmod {199}.
\]
By CRT this gives
\[
m\equiv 1\pmod {2197}.
\]

Thus I get the single residue \(m=1\) as the first possible value.  The probability is therefore
\[
\frac1{2197}.
\]
So \(p=1\), \(q=2197\), and
\[
p+q=2198\equiv \boxed{198}\pmod {1000}.
\]
